保留完整局部 top-$c$ 建立了正确性基础，但为了得到 $c$ 个全局输出，可能生成 $Gc$ 个中间候选。共享界组件 v63 减少这部分多余的选择工作，调度器 v65 再决定暴露多少扫描并行度。二者均在给定精排预算下工作；选择预算是\Cref{ch:refinement}讨论的独立问题。

\subsection{带明确见证的完整 Key 上界}

当 $|U|\geq c$ 时，用 $K_{(c)}$ 表示全局第 $c$ 小的 key。若某组持有 $c$ 个互不相同且已完整评分的候选，则这些见证中的最大 key $\tau_g$ 满足 $K_{(c)}\preceq\tau_g$：全局集合中已经存在 $c$ 个不晚于 $\tau_g$ 的候选。共享界初始化为最大哨兵，并按以下方式更新：
\begin{equation}
 \tau\leftarrow\min_{\prec}(\tau,\tau_g).
 \label{eq:5-shared-update}
\end{equation}
不足 $c$ 个有效候选时不得发布界。如果整个候选集合不足 $c$，任何组都无法发布这样的界，所有有效候选保持可接纳。

\begin{proposition}[由见证界安全排除候选]
\label{prop:5-shared-bound}
假设 key 固定，每次发布都有 $c$ 个互不相同的候选作为见证，每次读取返回初始哨兵或某个完整的已发布 key。仅在 $K_q(x)\succ\tau$ 时排除候选，不会删除任何全局 top-$c$ 成员，与发布时序无关。
\end{proposition}
\begin{proof}
每个已发布界都不小于 $K_{(c)}$，这些界的最小值也不小于 $K_{(c)}$。全局优胜候选满足 $K_q(x)\preceq K_{(c)}\preceq\tau$，因此不会被排除。较旧的界仍是上界，只会允许更多工作。此前的安全排除也不破坏后续发布的有效性：任意保留的 $c$ 个互不相同候选仍提供 $c$ 个见证，无论它们是否是该组最好的 $c$ 项。
\end{proof}

最后这一点很重要。启用共享后，一个组不必返回自己的完整局部 top-$c$：它可以丢掉局部有竞争力、但已被证明不可能进入全局结果的候选。正确性只要求各组输出的并集仍包含全部全局优胜候选，之后\Cref{prop:5-local-coverage}的截断合并论证仍适用。与共享界相等的候选继续有资格被接纳，包括建立该界的见证候选。实际接纳条件为
\begin{equation}
 K_q(x)\prec\lambda_g
 \quad\text{and}\quad
 K_q(x)\preceq\tau_{\rm seen},
 \label{eq:5-admission}
\end{equation}
其中 $\lambda_g$ 是局部压缩阈值，$\tau_{\rm seen}$ 是该组最近一次观察到的共享界。

\subsection{发布、过滤与实现边界}

对于由平方和生成的非负 FP32 分数，代码将距离位模式和非负位置打包为无符号 64-bit key：
\begin{equation}
 E(D,p)=2^{32}\operatorname{uintbits}_{32}(D)+p.
 \label{eq:5-packed-key}
\end{equation}
此时无符号整数比较与字典序一致。这依赖正零编码，并排除 NaN 和负分数；它不是通用的有符号浮点编码。正无穷排在有限分数之后，全一位模式的初始字又排在有效编码之后。无效位置不得用于发布界。这些是分数定义域与表示的前提，不能从 recall 不变推断出来。

\begin{figure}[tbp]
\centering
% Algorithm diagram; the whole-key-read premise is discussed in Section 5.5.
\begin{tikzpicture}[
  font=\small\linespread{1}\selectfont,
  box/.style={draw=ThesisNavy!65,rounded corners=2pt,align=center,
    inner sep=5pt,text width=6.2cm,minimum height=.9cm},
  side/.style={box,text width=4.1cm,fill=ThesisNavy!5},
  arr/.style={-{Latex[length=2mm]},draw=ThesisNavy!80,thick}]
\node[box] (score) at (0,0) {Read primary code; compute full $\primarydist$\\
  Form $K=(\primarydist,p)$};
\node[box,fill=ThesisNavy!5] (admit) at (0,-1.8)
  {Admission: $K\prec\lambda_j$ and $K\preceq\tau_{\rm seen}$};
\node[box] (compact) at (0,-3.8) {Append admitted keys\\
  Compact when necessary\\Keep at most $c$};
\node[box] (trim) at (0,-5.9) {Final compaction and prefix trim\\
  Emit valid keys $K\preceq\tau_{\rm seen}$ to merge};
\node[side] (shared) at (6.3,-.1) {Per-query shared word $Q$\\
  Encodes bound $\tau$\\Initially all ones\\$\tau\succeq K_{(c)}$};
\node[side] (publish) at (6.3,-3.8) {Only with $c$ valid witnesses:\\
  $\operatorname{atomicMin}(Q,\kappa(\tau_j))$};
\draw[arr] (score) -- (admit);
\draw[arr] (admit) -- node[right,font=\footnotesize] {yes} (compact);
\draw[arr] (compact) -- (trim);
\draw[arr] (compact.east) -- (publish.west);
\draw[arr] (publish.north) -- (shared.south);
\draw[arr,dashed] (shared.west) -- ++(-.35,0) |- (admit.east);
\node[side,minimum height=1.1cm] (discard) at (6.3,-6)
  {Discard candidate\\Scoring complete};
\draw[arr] (admit.west) -- node[above,font=\footnotesize] {no} ++(-.5,0) |- ([yshift=-.9cm]trim.south)
  -| (discard.south);
\end{tikzpicture}

\caption{单个扫描组的共享完整 key 过滤。候选的完整主码分数在接纳之前计算。有效的 $c$ 候选见证集合可以收紧查询级共享界；陈旧但完整的界只会削弱过滤。最终前缀裁剪减少合并的有效输入，但不改变已分配工作区的容量。完整 key 读取这一前提的实现边界在正文中另行说明。}
\label{fig:5-shared-bound}
\end{figure}

每个 block 的一个线程在每个 chunk 读取该查询的共享界。压缩后，如果有效条目足够，就通过 64-bit \texttt{atomicMin} 发布局部第 $c$ 小 key。扫描结束时，再进行一次发布，并通过二分查找只保留不晚于观察到的共享界的有序前缀。各 block 不进行跨 block 屏障同步，也不等待其他组。较晚的更新可能使另一组多输出一些条目，精确合并会安全删除它们。每次搜索前在同一 stream 中重置共享界，扫描 kernel 完成后才执行合并。

读取前提需要单独限定。归档代码以原子操作发布界，却通过 \texttt{volatile} 64-bit load 轮询。CUDA 不保证 \texttt{volatile} 访问具有原子性或线程间同步语义\cite{cuda_programming_guide}。要可移植地实现上述命题，读取端也应使用设备作用域的原子读取。\Cref{sec:7-4}中记录的 48 组配对检查支持被测构建中的实测一致性，不构成语言内存模型层面的并发证明。原子读取修订及其 GPU 复验尚未执行。

\Cref{fig:5-shared-bound}也明确了优化省下什么。每个被扫描候选仍读取主码，并计算完整主码分数。共享界避免一部分缓冲区追加、局部排序和中间写出，并缩短合并输入列表。它没有实现基于部分距离的提前终止，也没有减少完整主码评分的候选数量。收益还要抵消轮询和原子更新的开销；\Cref{sec:7-4}报告的配对加速范围中也包含轻微变慢的情况。

\subsection{在不改变预算的情况下选择并行度}

对 batch 大小 $b$，分组扫描启动 $bG$ 个 block，而不是 $b$ 个。这可以暴露原本不足的并行度，但也重复 LUT 搬运和局部选择，并增加合并步骤，因此固定组数的设计不可能在所有情形下都更好。较小静态配额可以减少中间存储，但要保持覆盖，需要额外证书或回退；pilot 观察到的优胜候选集中现象促使系统选择带见证的阈值方案（\Cref{sec:7-4}）。

实际调度器结合并行填充的近似指标和工作负载偏斜统计。令 $T$ 为扫描 block 的线程数，$H$ 为 SM 数乘以每个 SM 的最大驻留线程数，$v=\operatorname{sd}(n_c)/\operatorname{mean}(n_c)$ 为候选数的变异系数。规则为
\begin{align}
 G_{\rm fill}
 &=\max\!\left(\{1\}\cup
       \{g\in\{2,4,8\}:bgT\leq0.6H\}\right),
       \label{eq:5-fill-rule}\\
 G_{\rm skew}&=
   \begin{cases}4,&v\geq0.3,\\1,&v<0.3,\end{cases}
 \qquad
 G_{\rm chosen}=\max(G_{\rm fill},G_{\rm skew}).
 \label{eq:5-skew-rule}
\end{align}
偏斜规则可能覆盖填充限制。实验程序以 512 个数据库向量作为伪查询，在固定路由配置下校准 $v$，不使用 ground-truth 标签。这与第 6 章的采样输出一致性校准不同。校准之后，调度只依赖 batch 大小和已保存统计，不增加每次搜索的主机读回。两个常数在四个 pilot 配置上拟合，之后再到其他配置上评估。

当 $G_{\rm chosen}=1$ 时，实现选择更快的查询级路径，标签为 $G=0$。它不改变索引、探针数或精排预算，但继承该路径仅按距离截断时的并列行为。因此，在声称所有调度选择产生相同保留集合之前，必须检查共同全序前提。\Cref{sec:7-6}报告独立的集成输出检查，以及调度器的实测收益和 regret。

该规则是启发式策略，不是 occupancy 计算，也不是延迟 oracle：它没有建模寄存器压力、共享内存驻留限制、与预算相关的压缩成本和合并成本。冻结的伪查询统计也可能无法代表未来查询。记录中的最大 regret 出现在 batch 256，此时某个固定组数方案约快 $1.217\times$。因此，这里的设计贡献是经测试的候选并行扫描、带见证过滤与工作负载相关执行的组合，不是新的通用 top-$c$ 恒等式，也不保证最优调度。后续残差精排与最终 top-$k$ 按\Cref{sec:5-6}消费这些有效候选。
